\documentclass[11pt,a4paper]{article}
\usepackage[margin=28mm]{geometry}
\usepackage[T1]{fontenc}
\usepackage{lmodern,amsmath,amssymb,amsthm,booktabs,microtype}
\usepackage[hidelinks]{hyperref}
\usepackage{url}
\newtheorem{proposition}{Proposition}
\newtheorem{corollary}[proposition]{Corollary}
\newtheorem{definition}[proposition]{Definition}
\newcommand{\den}{\operatorname{den}}
\newcommand{\vp}{v_p}
\title{Prime neighbours of zeta denominators\\[4pt]\large Zeta doors: a computational note}
\author{Jack Pickett \textperiodcentered\ 22 September 2026}
\date{}
\begin{document}
\maketitle
\begin{abstract}
For positive indices $k\equiv2\pmod{12}$, let $D_k$ be the reduced positive
denominator of $\zeta(1-k)$. Classical Bernoulli arithmetic gives
$D_k=12T_k$, where every prime factor of $T_k$ is congruent to $11$ modulo
$12$. We select the unique neighbour of $D_k/3$ that is not divisible by
$3$ and investigate when it is prime. The construction provides known
factorisations on one side of each candidate and therefore supports
classical primality certificates. We give the elementary selection rule,
describe its limited Lean formalisation, and record computational examples
and the verification conditions needed for their interpretation. No claim
of an infinite prime family or of priority for the construction is made.
\end{abstract}

\section{The observation and the construction}
The equality
\[
 \zeta(-1)=\zeta(-13)=-\frac1{12}
\]
motivated an examination of arguments $1-k$ with $k\equiv2\pmod{12}$.
The equality itself does not establish a period of the zeta function.
Our object is the denominator at these arguments, rather than equality
of their values.

For a rational number $x$, write $\den(x)$ for its positive denominator
in lowest terms. For a prime $p$, let $v_p(k)$ denote the exponent of $p$
in the positive integer $k$. Throughout the note,
\[
 k\ge2,\qquad k\equiv2\pmod{12},\qquad
 D_k=\den\bigl(\zeta(1-k)\bigr).
\]
For an integer $n$ not divisible by $3$, put
\[
 \chi_3(n)=\begin{cases}1&n\equiv1\pmod3,\\-1&n\equiv2\pmod3.
 \end{cases}
\]
For an odd prime $p\ne3$, its \emph{3-free door} is
$m_0(p)=p+\chi_3(p)$: the even neighbour not divisible by $3$.

\section{Classical denominator arithmetic}
With the convention $B_1=-1/2$, the Bernoulli numbers satisfy
$\zeta(1-k)=-B_k/k$ for positive even $k$. The classical denominator
formula for these divided Bernoulli numbers is
\begin{equation}\label{eq:denominator}
 D_k=\prod_{\substack{p\text{ prime}\\p-1\mid k}}p^{1+v_p(k)}.
\end{equation}
For completeness, von Staudt--Clausen gives $v_p(B_k)=-1$ when
$p-1\mid k$. Hence $v_p(B_k/k)=-1-v_p(k)$ for those primes.
For the other primes, the classical integrality property of divided
Bernoulli numbers gives $B_k/k\in\mathbb Z_p$; see Kellner's account
\cite{kellner}. Together these facts give \eqref{eq:denominator}.
In particular, dividing the denominator of $B_k$ by $k$ without accounting
for cancellation would not justify this formula.

\begin{proposition}[Selection rule]\label{prop:selection}
If $p>3$ is prime and $p-1\mid k$, then
\[
 p\equiv11\pmod{12},\qquad m_0(p)=p-1\mid k.
\]
\end{proposition}
\begin{proof}
The index $k$ is even, is not divisible by $4$, and is not divisible by
$3$. Thus neither $4$ nor $3$ divides $p-1$. As $p$ is odd, the first
restriction forces $p\equiv3\pmod4$. Primality and $p>3$ exclude
$p\equiv0\pmod3$, and the second restriction excludes
$p\equiv1\pmod3$. Consequently $p\equiv2\pmod3$.
Combining these congruences gives $p\equiv11\pmod{12}$.
The definition of $m_0$ then gives $m_0(p)=p-1$.
\end{proof}

\begin{corollary}\label{cor:shape}
There is an odd positive integer $T_k$ such that
\[
 D_k=12T_k,\qquad
 T_k=\prod_{\substack{p>3\text{ prime}\\p-1\mid k}}p^{1+v_p(k)}.
\]
Every prime factor of $T_k$ is $11$ modulo $12$. In particular,
$d_k=D_k/3=4T_k$ is even and not divisible by $3$.
\end{corollary}
\begin{proof}
In \eqref{eq:denominator}, $2$ occurs to exponent $2$, since $v_2(k)=1$,
and $3$ occurs to exponent $1$, since $v_3(k)=0$. Apply
Proposition~\ref{prop:selection} to the remaining factors.
\end{proof}

\begin{definition}
The \emph{zeta-door candidate} at $k$ is
\[
 Q_k=d_k+\chi_3(d_k).
\]
We call it a zeta-door prime when $Q_k$ is prime.
\end{definition}
This selects exactly one of $d_k-1$ and $d_k+1$: the other is divisible
by $3$. If $Q_k$ is prime, its residue modulo $3$ is opposite to that
of $d_k$, so
\begin{equation}\label{eq:door}
 m_0(Q_k)=d_k.
\end{equation}
Writing $\Omega(T)$ for the number of prime factors counted with
multiplicity, including $\Omega(1)=0$, an equivalent expression is
\[
 Q_k=4T_k+(-1)^{\Omega(T_k)}.
\]
The exponent multiplicities in this expression are essential.

\section{Small examples and the relation to doors}
Direct rational evaluation gives the following examples.
\begin{center}
\begin{tabular}{rrrrr}
\toprule
$k$ & $1-k$ & $D_k$ & $d_k$ & $Q_k$\\
\midrule
2 & $-1$ & 12 & 4 & 5\\
14 & $-13$ & 12 & 4 & 5\\
26 & $-25$ & 12 & 4 & 5\\
38 & $-37$ & 12 & 4 & 5\\
50 & $-49$ & 132 & 44 & 43\\
\bottomrule
\end{tabular}
\end{center}
Thus different indices can produce the same candidate. Counts of
indices, distinct denominators, and distinct primes must be separated.
The table concerns reduced denominators; it does not assert that the
zeta values in its first four rows are equal.

Two divisibility relations also require separation. The denominator
selection rule reads $m_0(p)\mid k$ for a contributing prime $p>3$.
The ownership relation in the hire construction reads $p\mid m_0(Q)$.
When $Q=Q_k$ is prime, \eqref{eq:door} implies that every prime factor
of $T_k$ is an owner-door factor of $Q$. These observations do not
identify the index $k$ with a graph vertex or establish a theorem about
finite gold components or Laplacian spectra.

\section{Search and certification}
The search starts with a fully factored index $k$. It enumerates all
positive divisors $d\mid k$ and examines $p=d+1$ for primality.
Equation~\eqref{eq:denominator} then constructs $D_k$ without computing
the enormous numerator of $B_k/k$. Small-factor sieving and probable-prime
tests screen the resulting $Q_k$. Screening is distinct from proof.

The targeted runs recorded here use the plus branch $Q_k=d_k+1$.
Consequently $Q_k-1=d_k$ has a supplied factorisation. If $N=Q_k$ and
$N-1=FR$, where $F$ is completely factored into proved primes and
$F^2>N$, a partial Pocklington certificate suffices. For each prime
$r\mid F$, the certificate supplies a witness $a_r$ with
\[
 a_r^{N-1}\equiv1\pmod N,\qquad
 \gcd\bigl(a_r^{(N-1)/r}-1,N\bigr)=1.
\]
These conditions force every prime divisor of $N$ to be $1$ modulo $F$;
since $F>\sqrt N$, $N$ must be prime. This standard criterion is the
primality engine, not a new primality theorem; see \cite{pocklington}.
Minus-branch candidates would instead expose a factorisation of $Q_k+1$
and need a suitable different certification method.

Two independent obligations remain:
\begin{enumerate}
\item Prove the candidate $Q_k$ prime. Only enough proved factors of
$Q_k-1$ to meet Pocklington's bound are needed.
\item Prove the claimed denominator is exact. Every included $d+1$ must
be prime, and every omitted $d+1$ must be composite.
\end{enumerate}
For a large contributing prime $p$, the factorisation of $p-1\mid k$
is already available from that of $k$. A complete $p-1$ certificate can
therefore establish its primality at relatively low cost. The verifier
uses a deterministic Miller--Rabin test in the range $n<2^{64}$ and
Pocklington above that range. Its seven Miller--Rabin bases are
\[
 2,\ 325,\ 9375,\ 28178,\ 450775,\ 9780504,\ 1795265022.
\]
The range restriction is part of the verification contract.

\section{Computational snapshot and reproducibility}
The notebook reports the following successful candidate-primality
checks in the \texttt{kitchen8} run. These are experimental records of
this search, not claims of global size records.
\begin{center}
\begin{tabular}{rrr}
\toprule
Digits of $Q_k$ & Distinct primes below $2^{64}$ in $F$
 & Larger listed factors\\
\midrule
5522 & 510 & 5\\
5118 & 462 & 3\\
4443 & 360 & 7\\
\bottomrule
\end{tabular}
\end{center}
For the 5522-digit entry the index is
\[
 k=111695474196999960960050.
\]
The author reports that all $5$, $3$, and $7$ larger listed factors,
respectively, were subsequently proved prime from their index
factorisations. A separate 18365-digit candidate at
$k=1522610689958516981964050$ remains a probable prime in this snapshot;
its certificate computation is not included among completed results.

The source snapshot for this draft is repository revision
\texttt{8b99e2b} \cite{repo}. Change into the directory
\begin{center}
\path{analysis/zeta/zeta-targeted}
\end{center}
and run the main reproducibility command:
\begin{verbatim}
python3 verify.py kitchen8/hits.jsonl
\end{verbatim}
Use ordinary Python execution, without optimisation flags that disable
assertions. The JSON lines contain the exact candidates, indices,
factorisations, and candidate witnesses; they are preferable to printing
thousands of decimal digits in the note. Reported long certificate runs
have not been independently repeated during preparation of this draft.

\paragraph{Outstanding audit item for this draft.}
Revision \texttt{8b99e2b} searches bases $2,\ldots,501$ when proving
large divisor-derived candidates. Exhausting this bounded search returns
\texttt{False}, which the denominator loop treats as a reason to omit the
candidate. Failure to find a primality witness is not a compositeness
proof. A publication-ready exact-origin receipt must distinguish
\emph{prime}, \emph{composite}, and \emph{unresolved}, and must not finish
successfully with unresolved divisor candidates. The successful proofs
of all included large factors do not by themselves settle this
exhaustiveness issue. This qualification concerns denominator provenance,
not the separate partial-Pocklington proof of $Q_k$.

\section{Formalisation and further questions}
The file \texttt{lean/Hire/ZetaDoors.lean} contains the selection theorem
\begin{center}
\texttt{prime\_mod\_twelve\_of\_pred\_dvd\_index}
\end{center}
and its door-divisibility consequence
\texttt{m0\_dvd\_of\_pred\_dvd\_index}. The first proves
Proposition~\ref{prop:selection}'s residue conclusion from primality,
$p\ne2,3$, $k\bmod12=2$, and $p-1\mid k$; the second rewrites the
door as $p-1$. The file mentions Mathlib's von Staudt--Clausen theorem
but does not invoke it. It does not formalise
\eqref{eq:denominator}, the complete search, or the large numerical
primality certificates.

The immediate computational questions are how often denominators repeat,
how the two signs are distributed, and how prime frequency changes under
different index-selection schemes. A density measured in a targeted
sample of highly factored indices should not be presented as the density
over all $k\equiv2\pmod{12}$. Whether infinitely many distinct zeta-door
primes exist is not settled here. A literature comparison is also needed
before making a priority claim for this particular construction.

The companion Lab at \url{https://halfasecond.com/} retains changing
records and experiments; this note fixes definitions, proof scope, and a
dated computational snapshot.

\begin{thebibliography}{9}
\bibitem{kellner} B. C. Kellner, \emph{The Bernoulli Number Page},
sections on denominator and numerator structure,
\url{https://www.bernoulli.org/}. Accessed 22 September 2026.
\bibitem{pocklington} PrimePages, \emph{Finding primes and proving
primality}, Section 3: classical primality tests,
\url{https://t5k.org/prove/prove3.html}.
\bibitem{repo} J. Pickett, \emph{Natural Maths notebook}, source and
computational records, revision \texttt{8b99e2b},
\url{https://github.com/hasjack/nm-notebook/tree/8b99e2b}.
\end{thebibliography}
\end{document}
